% ==============================================================================
\section{Python Reference Implementation}
\label{sec:python-implementation}
% ==============================================================================

We include a simplified Python implementation of the Fast CASPaxos acceptor and proposer as a pedagogical complement to Figure~\ref{fig:fast-caspaxos-summary}.
A fuller demo harness is available in the accompanying repository.

\paragraph{Ballots and quorums.}
A ballot is a \texttt{(round, proposer)} pair ordered lexicographically; \texttt{proposer\,=\,0} denotes a fast ballot.
Classic quorums require $\lfloor N/2 \rfloor + 1$ acceptors; fast quorums require $\lceil 3N/4 \rceil$.

\paragraph{Acceptor.}
The acceptor handles \textsc{Prepare} and \textsc{Accept} RPCs.
Its state matches Figure~\ref{fig:fast-caspaxos-summary}: \texttt{promised}, \texttt{accepted\_ballot}, and \texttt{accepted\_value}.
The \textsc{Accept} handler uses first-write-wins for fast ballots and processes the piggybacked prepare atomically.

\begin{lstlisting}[
  style=pseudocode,
  numbers=left,
  numberstyle=\tiny\color{gray},
  xleftmargin=2em,
  basicstyle=\scriptsize\ttfamily,
  aboveskip=6pt,
  belowskip=6pt,
]
class Acceptor:
    def __init__(self, name, initial_value=None):
        self.promised = Ballot(1, 0)  # initial fast ballot
        self.accepted_ballot = Ballot(0, 0)
        self.accepted_value = initial_value

    def prepare(self, ballot):                    # Prepare RPC
        if ballot < max(self.promised, self.accepted_ballot):
            return Reject(self.promised)
        self.promised = ballot
        return Promise(self.accepted_ballot, self.accepted_value)

    def accept(self, ballot, value, next_ballot=None):
        # Idempotent replay (same ballot and value)
        if ballot == self.accepted_ballot \
                and value == self.accepted_value:
            return Accepted(self.promised)
        if ballot < max(self.promised, self.accepted_ballot):
            return Reject(self.promised)
        # Same ballot, different value: reject
        if ballot == self.accepted_ballot:
            return Reject(self.promised)
        self.accepted_ballot = ballot
        self.accepted_value = value
        # Piggybacked prepare: prepare next round
        if next_ballot and next_ballot > self.promised:
            self.promised = next_ballot
        return Accepted(self.promised)
\end{lstlisting}

\paragraph{Proposer.}
The proposer runs the Read--Modify--Write cycle.
It reuses a piggybacked prepare when available; otherwise it runs a full \textsc{Prepare}.
The update function is applied locally, and the result is committed with \textsc{Accept}.
With full responses, fast-ballot recovery in \texttt{\_recover} reduces to choosing the unique maximum, or any tied maximum, at the highest fast ballot observed.

\paragraph{Learning.}
The reference implementation models inter-proposer learning as a best-effort notification carrying the decided ballot, value, and prepared next ballot.
If that notification is missing, the proposer falls back to \textsc{Prepare} on its next operation.

\paragraph{Simplifications.}
This reference implementation assumes that all acceptors respond; a production implementation must handle partial responses and retry with a higher classic ballot when fast-ballot recovery remains ambiguous.

\begin{lstlisting}[
  style=pseudocode,
  numbers=left,
  numberstyle=\tiny\color{gray},
  xleftmargin=2em,
  basicstyle=\scriptsize\ttfamily,
  aboveskip=6pt,
  belowskip=6pt,
]
class Proposer:
    def __init__(self, id, acceptors, enable_fast=True):
        self.id = id
        self.acceptors = acceptors
        self.enable_fast = enable_fast
        self.ballot = Ballot(0, 0)
        self.cached_value = None
        self.prepared = None  # piggybacked prepare ballot

    def propose(self, update_fn):
        while True:
            # 1. Read: obtain ballot and current value
            b, need_read = self._choose_ballot()
            if need_read:
                ok, val = self._run_read(b)
                if not ok: continue
                self.cached_value = val
            # 2. Modify: apply the client's update function
            new_val = update_fn(self.cached_value)
            # 3. Write: commit the new value
            if self._run_write(b, new_val):
                self.cached_value = new_val
                return new_val

    def _choose_ballot(self):
        if self.prepared and self.prepared > self.ballot:
            b, self.prepared = self.prepared, None
            self.ballot = b
            return b, False         # reuse piggybacked prepare
        if self.enable_fast \
                and self.ballot < Ballot(1, 0):  # initial fast
            self.ballot = Ballot(1, 0)
            return self.ballot, False  # initial fast ballot
        self.ballot = Ballot(self.ballot.round + 1, self.id)
        return self.ballot, True    # classic: read required

    def _run_read(self, ballot):
        resps = [a.prepare(ballot) for a in self.acceptors]
        promises = [r for r in resps if isinstance(r, Promise)]
        N = len(self.acceptors)
        if len(promises) < N // 2 + 1:
            return False, None
        return True, self._recover(promises, N)

    def _recover(self, promises, N):
        best = max(promises, key=lambda p: p.accepted_ballot)
        if best.accepted_ballot == Ballot(0, 0):
            return self.cached_value
        if best.accepted_ballot.proposer != 0:  # classic
            return best.accepted_value
        # Fast ballot: choose any tied maximum at best ballot
        at_best = [p.accepted_value for p in promises
                   if p.accepted_ballot == best.accepted_ballot]
        counts = Counter(at_best)
        max_count = max(counts.values())
        tied = [v for v, c in counts.items() if c == max_count]
        return tied[0]

    def _run_write(self, ballot, value):
        N = len(self.acceptors)
        next_b = Ballot(ballot.round + 1,
                        0 if self.enable_fast else self.id)
        resps = [a.accept(ballot, value, next_b)
                 for a in self.acceptors]
        accepts = [r for r in resps if isinstance(r, Accepted)]
        needed = -(-3*N//4) if ballot.proposer == 0 \
                 else N//2 + 1
        if len(accepts) >= needed:
            self.prepared = next_b  # cache piggybacked prepare
            return True
        return False
\end{lstlisting}
